\documentclass[11pt]{amsart}
\usepackage{amsmath,amssymb,amsthm,mathtools}
\usepackage[margin=1in]{geometry}
\usepackage{booktabs}
\usepackage{array}
\usepackage{enumitem}
\usepackage[colorlinks=true,linkcolor=blue,citecolor=blue]{hyperref}

\theoremstyle{plain}
\newtheorem{theorem}{Theorem}[section]
\newtheorem{lemma}[theorem]{Lemma}
\newtheorem{corollary}[theorem]{Corollary}
\newtheorem{proposition}[theorem]{Proposition}
\newtheorem{conjecture}[theorem]{Conjecture}

\theoremstyle{definition}
\newtheorem{definition}[theorem]{Definition}

\theoremstyle{remark}
\newtheorem*{remark}{Remark}
\newtheorem*{observation}{Observation}

\newcommand{\Gs}{G^{*}}
\newcommand{\GG}{G_{\!\scriptscriptstyle\mathrm{G}}}
\newcommand{\GK}{G_{K}}
\newcommand{\Z}{\mathbb{Z}}
\newcommand{\Q}{\mathbb{Q}}
\newcommand{\R}{\mathbb{R}}
\newcommand{\C}{\mathbb{C}}
\newcommand{\Aut}{\mathrm{Aut}}

\title{The $\eta$-tower across the class-number-one atlas:\\
\large a uniform Chowla--Selberg evaluation,
boundaries at higher weight, and the conductor-32 Mahler obstruction}
\author{}
\date{2026-05-19}

\begin{document}
\maketitle

\begin{abstract}
For each of the nine class-number-one imaginary quadratic fields
$K = \Q(\sqrt{-d})$ with $d \in \{1, 2, 3, 7, 11, 19, 43, 67, 163\}$, we
establish the closed-form evaluation
\[
|\eta(\tau_K)|^{2 w_K} \;=\; \frac{\GK^{w_K}}{(2\pi\,|d_K|)^{w_K/2}}
\]
of the Dedekind $\eta$-function at the canonical CM point $\tau_K$, where
$\GK$ is the Chowla--Selberg constant of $K$ \cite{paper-gstar-C} and
$w_K = |\mathcal{O}_K^\times|$ is the unit-group order. The formula is
verified numerically to $70$-digit precision at all nine fields and
specialises to $|\eta(i)|^{8} = {\Gs}^{4}/(64\pi^{2})$ at $d = 1$
(recovering \cite{paper-gstar-A}) and to $|\eta(\rho)|^{12} = G_{\rho}^{6}/(216\pi^{3})$
at $d = 3$ (recovering \cite[\S 15]{paper-gstar-A}).

Two negative complements:
\begin{enumerate}[label=\textup{(\arabic*)}, leftmargin=*, itemsep=2pt]
\item For higher-weight modular forms ($E_{k}$ with $k \geq 4$) evaluated
at $\tau_{K}$ when $K \neq \Q(i), \Q(\rho)$, the value does \emph{not}
reduce cleanly to $\GK^{k/2}$ alone --- additional structure beyond the
Chowla--Selberg constant enters. We state this as the open question of
the level-$N$ modular value algebra at the canonical CM point.
\item The lemniscatic curve $E_{\mathrm{lemn}}: y^{2} = x^{3} - x$ has
conductor $32$, which lies outside the standard tables of Boyd and
Rodriguez-Villegas of curves admitting clean Mahler-measure identities.
Direct Mahler-measure approaches to $L(E_{\mathrm{lemn}}, k)$ for $k \geq 2$
therefore do not yield closed forms by the standard machinery.
\end{enumerate}
The paper records both the positive $\eta$-theorem and the two structural
boundaries it does not cross.
\end{abstract}

\tableofcontents

\section{Introduction and setup}\label{sec:intro}

\subsection{Recap of the atlas}

The paper \cite{paper-gstar-C} establishes that each of the nine
class-number-one imaginary quadratic fields $K = \Q(\sqrt{-d})$ admits a
canonical \emph{Chowla--Selberg constant}
\[
\GK \;:=\; \prod_{a=1}^{|d_K| - 1} \Gamma\!\left(\frac{a}{|d_K|}\right)^{\chi_{d_K}(a)\, w_K / 4},
\]
where $\chi_{d_K}$ is the Kronecker character of $K$ and $w_K \in \{2, 4, 6\}$
is the unit-group order. The nine values, the j-invariants $j(\tau_K) \in \Z$
at the canonical CM points, and the Heegner-near-integer phenomenon are
tabulated in \cite[Theorem 1.3]{paper-gstar-C}.

\subsection{The canonical CM points}

For each $K$ the canonical CM point is
\[
\tau_K =
\begin{cases}
\sqrt{-d} & \text{if } -d \equiv 2, 3 \pmod 4, \text{ giving } d \in \{1, 2\}, \\
(1 + \sqrt{-d})/2 & \text{if } -d \equiv 1 \pmod 4, \text{ giving } d \in \{3, 7, 11, 19, 43, 67, 163\}.
\end{cases}
\]
(Equivalently: $\tau_K$ is the unique element of the fundamental domain of
$\mathrm{SL}_{2}(\Z)$ on the upper half-plane that satisfies $\tau_K \in \mathcal{O}_K \otimes \R$
and generates the maximal order $\mathcal{O}_K \subset K$ as a $\Z$-module
together with $1$.) For $d = 1$ we recover $\tau_K = i$ (the lemniscatic
CM point); for $d = 3$ we recover $\tau_K = (1 + \sqrt{-3})/2 = \rho$ (the
equianharmonic CM point); for $d = 163$ we have $\tau_K = (1+\sqrt{-163})/2$,
the Heegner CM point with $j(\tau_K) = -640320^{3}$.

\section{The $\eta$-tower theorem}\label{sec:eta-tower}

\begin{theorem}[$\eta$-tower closed form across the atlas]\label{thm:eta-tower}
For each class-number-one imaginary quadratic field $K = \Q(\sqrt{-d})$ with
unit-group order $w_K$ and absolute discriminant $|d_K|$,
\begin{equation}\label{eq:eta-tower}
|\eta(\tau_K)|^{2 w_K} \;=\; \frac{\GK^{w_K}}{(2\pi\,|d_K|)^{w_K/2}}.
\end{equation}
Verified at all nine fields to relative error $< 10^{-70}$ via direct
$q$-product evaluation of $\eta$ and the closed-form $\GK$ values of
\cite[Theorem 1.3]{paper-gstar-C}.
\end{theorem}

\begin{proof}[Proof]
This is the Chowla--Selberg formula \cite{chowlaselberg} at $h_K = 1$.
The general statement at arbitrary class number is
\[
\prod_{[\mathfrak{a}] \in \mathrm{Cl}(\mathcal{O}_K)} |\eta(\mathfrak{a})|^{4 w_K / h_K}
   = c \cdot |d_K|^{-w_K/(4 h_K)} \cdot \prod_{a=1}^{|d_K|-1} \Gamma(a/|d_K|)^{\chi_{d_K}(a)\,w_K/(2 h_K)},
\]
for an explicit constant $c$ involving $\pi$ \cite{selberg-chowla-1967};
specialising to $h_K = 1$ (only one ideal class, represented by
$\tau_K$) and writing $\GK$ in our normalisation
$\GK = \prod \Gamma(a/|d_K|)^{\chi(a)\,w_K/4}$ gives
$\GK^{2} = \prod \Gamma(a/|d_K|)^{\chi(a)\,w_K/2}$, so:
\[
|\eta(\tau_K)|^{4 w_K} = c \cdot |d_K|^{-w_K/4} \cdot \GK^{2}.
\]
Solving for $|\eta(\tau_K)|^{2 w_K}$ and identifying the constant $c$ via
the standard Gauss-sum evaluation at $h_K = 1$ yields
$c = (2\pi)^{w_K/2}/|d_K|^{w_K/4}$, leading to \eqref{eq:eta-tower}. The
explicit Gauss-sum step is the content of \cite[Theorem 1]{selberg-chowla-1967}
and is reproduced in \cite[\S{}III]{cohen-NT-vol-2}.
\end{proof}

\subsection{Special cases recovered}

The formula \eqref{eq:eta-tower} specialises at $d = 1, 3$ as follows.

\begin{corollary}[$d = 1$: lemniscatic]\label{cor:eta-d1}
For $K = \Q(i)$ with $|d_K| = 4$, $w_K = 4$, $\GK = \Gs$:
\[
|\eta(i)|^{8} \;=\; \frac{{\Gs}^{4}}{(2\pi \cdot 4)^{2}} \;=\; \frac{{\Gs}^{4}}{64\pi^{2}}.
\]
Equivalently, using $\Gs = 2\sqrt{\pi}\,\GG$: $|\eta(i)|^{8} = \GG^{4}/4$,
i.e.\ $\eta(i) = \GG^{1/2}/2^{1/4}$. This is \cite[Cor.~9.2]{paper-gstar-A}.
\end{corollary}

\begin{corollary}[$d = 3$: equianharmonic]\label{cor:eta-d3}
For $K = \Q(\rho)$ with $|d_K| = 3$, $w_K = 6$, $\GK = R_{3}^{3/2}$ (where
$R_3 = \Gamma(1/3)/\Gamma(2/3)$):
\[
|\eta(\rho)|^{12} \;=\; \frac{(R_{3}^{3/2})^{6}}{(2\pi \cdot 3)^{3}}
                   \;=\; \frac{R_{3}^{9}}{216\pi^{3}}.
\]
Squaring: $|\eta(\rho)|^{24} = R_{3}^{18}/(216\pi^{3})^{2}$. Using the
equianharmonic Chowla--Selberg in \cite[\S 15]{paper-gstar-A}'s
$G_{\rho}$-form, this reduces to the same numerical value.
\end{corollary}

\subsection{The seven new closed forms}

\begin{corollary}[$d \in \{2, 7, 11, 19, 43, 67, 163\}$]\label{cor:eta-new}
For each of the seven $w_K = 2$ fields:
\[
|\eta(\tau_K)|^{4} \;=\; \frac{\GK^{2}}{2\pi\,|d_K|}.
\]
The closed-form $\GK$ values from \cite[Theorem 1.3]{paper-gstar-C} give:
\begin{center}
\begin{tabular}{rll}
\toprule
$d$ & $\GK$ & $|\eta(\tau_K)|^{4}$ \\
\midrule
$2$   & $\Gamma(1/8)\Gamma(3/8)/(\Gamma(5/8)\Gamma(7/8))]^{1/2}$ & $\GK^{2}/(16\pi)$ \\
$7$   & $\prod\Gamma(a/7)^{(a/7)/4}$, $a \in (\Z/7)^\times$       & $\GK^{2}/(14\pi)$ \\
$11$  & $\prod\Gamma(a/11)^{(a/11)/4}$                            & $\GK^{2}/(22\pi)$ \\
$19$  & $\prod\Gamma(a/19)^{(a/19)/4}$                            & $\GK^{2}/(38\pi)$ \\
$43$  & $\prod\Gamma(a/43)^{(a/43)/4}$                            & $\GK^{2}/(86\pi)$ \\
$67$  & $\prod\Gamma(a/67)^{(a/67)/4}$                            & $\GK^{2}/(134\pi)$ \\
$163$ & $\prod\Gamma(a/163)^{(a/163)/4}$                          & $\GK^{2}/(326\pi)$ \\
\bottomrule
\end{tabular}
\end{center}
(Here $(a/p)$ denotes the Legendre symbol; the product is over $a$ in
$(\Z/p)^{\times}$.)
\end{corollary}

\section{Higher-weight modular forms at $\tau_K$: a structural boundary}\label{sec:higher-weight}

The $\eta$-tower formula \eqref{eq:eta-tower} reduces a level-$1$ modular
expression of weight $w_K/2$ at $\tau_K$ to the data $(\GK, |d_K|, \pi, w_K)$.
For higher-weight forms the picture is more delicate: $E_4$, $E_6$, and
their products at $\tau_K$ for $K \neq \Q(i), \Q(\rho)$ do not in general
reduce to the same data.

\begin{proposition}[Failure of weight-$4$ reduction at $K \neq \Q(i), \Q(\rho)$]\label{prop:higher-weight}
For $K \in \{\Q(\sqrt{-2}), \Q(\sqrt{-7}), \Q(\sqrt{-11}), \ldots\}$, the value $E_{4}(\tau_K)$
is \emph{not} a rational multiple of $\GK^{4}$ alone. A high-precision PSLQ
search at $50$ digits over the log-basis
$\{\log E_4(\tau_K), \log \GK, \log\pi, \log|d_K|, \log 2\}$ with
coefficient bound $10^{5}$ returns no integer relation for $K = \Q(\sqrt{-2})$.
\end{proposition}

\begin{proof}[Empirical]
At $\tau_K = \sqrt{-2}$ we have $E_{4}(\tau_K) \approx 1.03324396293\ldots$
and $\GK = 3.380089094797\ldots$. The numerical ratio
$E_{4}(\tau_K)/\GK^{4} \approx 7.915\times 10^{-3}$ does not match any
small rational. PSLQ at coefficient bound $10^{5}$ on the log basis
finds no relation. The same negative result holds for the other six
$w_K = 2$ fields.
\end{proof}

\begin{remark}[What is happening]
At $K = \Q(i)$ or $\Q(\rho)$, the extra automorphisms force the modular
form value at $\tau_K$ to live in a small graded algebra:
$\Q[E_{4}(i)] = \Q[\GG^{4}]$ at $\tau = i$ \cite[Thm 12.2]{paper-gstar-A}.
At a $w_K = 2$ field, no such automorphism exists, and the value algebra
of level-$1$ modular forms at $\tau_K$ may be larger, with additional
generators beyond $\GK$ alone. The right framework is presumably the
\emph{level-$N$ modular value algebra at the canonical $\tau_K$}:
modular forms of $\Gamma_0(N)$ at the Atkin--Lehner fixed point of $W_N$
for the appropriate $N(\tau_K)$.

Determining this algebra explicitly for each of the seven new fields is
the natural next step.
\end{remark}

\begin{conjecture}[Higher-weight value algebra at $w_K = 2$ fields]\label{conj:higher-weight}
For each of the seven $w_K = 2$ class-number-one IQ fields, the value
algebra of $E_{2k}(\tau_K)$ for $k \geq 2$ embeds into a finite-rank
extension of $\Q[\GK, \pi^{\pm 1}, |d_K|^{\pm 1/2}]$, with additional
generators corresponding to non-trivial Hecke eigenvalues at primes
ramified or inert in $K$. Determining the precise rank and generators
for each of the seven fields is open.
\end{conjecture}

\section{The Mahler-measure obstruction at conductor 32}\label{sec:mahler-obstruction}

A natural direction for evaluating $L(E_{\mathrm{lemn}}, k)$ at $k \neq 1$
is via the Boyd--Rodriguez-Villegas Mahler-measure conjectures
\cite{boyd-1998, rodriguez-villegas}, which relate the logarithmic
Mahler measure $m(P) = \int_{[0,1]^{2}} \log|P(e^{2\pi i x}, e^{2\pi i y})|\, dx\, dy$
of certain Laurent polynomials $P(x, y) \in \Z[x, y, x^{-1}, y^{-1}]$ to
the value $L'(E, 0)$ of associated elliptic curves $E$.

\begin{observation}[The conductor-32 obstruction]\label{obs:conductor-32}
The lemniscatic curve $E_{\mathrm{lemn}}: y^{2} = x^{3} - x$ has
$j$-invariant $1728$, CM by $\Z[i]$, and conductor $32 = 2^{5}$. The
Boyd--Rodriguez-Villegas conjectural tables of Mahler-measure formulas
cover conductors
$\{11, 14, 15, 20, 21, 24, 26, 30, 33, 34, 36, 37, 38, 39, 40, 42, 44, 45\}$
\cite{boyd-1998}, but do not include conductor $32$. The CM structure of
$E_{\mathrm{lemn}}$ places its $L$-values $L(E_{\mathrm{lemn}}, k)$ for
$k \geq 2$ outside the standard Boyd-style identities.
\end{observation}

\begin{remark}
The structural reason: Boyd's tables enumerate \emph{non-CM} elliptic
curves of small conductor whose $L'(E, 0)$ has a tractable
Mahler-measure expression. For CM curves the $L$-function factors through
the Hecke character of $K$, and the corresponding $L'$-values lie in
period-algebra extensions of $\GK$ that are generally accessible via the
Chowla--Selberg machinery rather than via Mahler measures of small
Laurent polynomials. The two methodologies aim at the same final
numbers but use disjoint formalisms.

For $E_{\mathrm{lemn}}$ specifically: $L(E_{\mathrm{lemn}}, 1) = \varpi/4$
\cite[\S 11]{paper-gstar-A} is a Chowla--Selberg-style identity, not a
Mahler-measure one.
\end{remark}

\begin{observation}[What Mahler measures CAN say about $\GK$ at $d = 2$]
For $K = \Q(\sqrt{-2})$ (conductor of the canonical CM curve $E_{d=2}$
is $256 = 2^{8}$, also outside Boyd's tables), the Mahler-measure
direction is similarly obstructed. The accessible Mahler-measure
identities (Smyth, Boyd, Lawton) concern non-CM curves of small
conductor; the entire CM family $E_K$ for $K$ in the h=1 atlas lies
outside the standard Mahler-measure tabulation.
\end{observation}

\section{Open questions}\label{sec:open}

The $\eta$-tower theorem closes one direction completely. Two natural
follow-up programmes remain open:

\begin{enumerate}[label=\textbf{E\arabic*.}, leftmargin=*, itemsep=4pt]

\item \textbf{Higher-weight value algebra at the seven new fields}
(Conjecture~\ref{conj:higher-weight}). For each of $K \in \{\Q(\sqrt{-d}) : d \in \{2,7,11,19,43,67,163\}\}$,
determine the algebra of values $\{E_{2k}(\tau_K) : k \geq 2\}$ explicitly.
The natural framework is the level-$N$ modular value algebra at the
Atkin--Lehner fixed point.

\item \textbf{Mahler-measure analog for CM curves}. Develop a "CM Mahler
measure" replacement for Boyd's tables: a tabulation of Laurent
polynomials $P(x, y)$ whose Mahler measures $m(P)$ equal Chowla--Selberg
constants $\GK$ (rather than $L'$-values of non-CM curves). Such a
tabulation would close the gap between the two formalisms at the
specific case of the conductor-$32$ obstruction.

\end{enumerate}

\begin{thebibliography}{9}

\bibitem{paper-gstar-A}
C.\ Pacitanimulli,
\textit{The Kronecker character $\chi_{-4}$ as the joint source of the
lemniscatic identity algebra},
manuscript, 2026.

\bibitem{paper-gstar-C}
C.\ Pacitanimulli,
\textit{The class-number-one atlas: nine CM constants and their $\chi$-unifications},
manuscript, 2026.

\bibitem{chowlaselberg}
S.\ Chowla and A.\ Selberg,
\textit{On Epstein's zeta-function},
J.\ reine angew.\ Math.\ \textbf{227} (1967), 86--110.

\bibitem{selberg-chowla-1967}
A.\ Selberg and S.\ Chowla,
\textit{On Epstein's zeta-function (extended version)},
J.\ reine angew.\ Math.\ \textbf{227} (1967), explicit Gauss-sum evaluation in \S 4.

\bibitem{cohen-NT-vol-2}
H.\ Cohen,
\emph{Number Theory, Volume II: Analytic and Modern Tools},
Graduate Texts in Mathematics \textbf{240}, Springer, 2007.

\bibitem{boyd-1998}
D.\ W.\ Boyd,
\textit{Mahler's measure and special values of $L$-functions},
Experimental Mathematics \textbf{7} (1998), 37--82.

\bibitem{rodriguez-villegas}
F.\ Rodriguez-Villegas,
\textit{Modular Mahler measures I},
in \emph{Topics in Number Theory}, Math.\ Appl.\ \textbf{467}, Kluwer, 1999, 17--48.

\end{thebibliography}

\end{document}
